\documentclass[a4paper,12pt]{article}
\usepackage[utf8]{inputenc}
\usepackage[T1]{fontenc}
\usepackage[ngerman]{babel}
\usepackage{geometry}
\geometry{margin=2.5cm}
\usepackage{amsmath}
\usepackage{booktabs}
\usepackage{hyperref}
\hypersetup{colorlinks=true, linkcolor=blue, urlcolor=blue}

\begin{document}

% =========================================================
\section*{Convolutional Neural Operators for Robust and Accurate
          Learning of PDEs}
% =========================================================

\noindent\textbf{Autoren:} Bogdan Raoni\'c, Roberto Molinaro, Tim De Ryck,
Tobias Rohner, Francesca Bartolucci, Rima Alaifari, Siddhartha Mishra,
Emmanuel de B\'ezenac\\[2pt]
\textbf{Institution:} Seminar for Applied Mathematics, ETH Z\"urich;
Delft University of Technology\\[2pt]
\textbf{Ver\"offentlicht:} NeurIPS 2023 (arXiv:2302.01178v3, Dezember 2023)\\[2pt]
\textbf{Code:} \url{https://github.com/bogdanraonic3/ConvolutionalNeuralOperator}

% =========================================================
\subsection*{Motivation und Problemstellung}
% =========================================================

Klassische faltungsbasierte Netzwerke (CNNs) gelten im Kontext des
Operator"=Lernens bislang als konzeptionell problematisch, da ihre naive
Anwendung zu sogenannten Aliasing-Fehlern f\"uhrt: Die diskrete
Implementierung eines CNNs entspricht nicht konsistent einer kontinuierlichen
Abbildung zwischen Funktionenr\"aumen. Dies steht im Widerspruch zur
Kernaufgabe des Operator"=Lernens, n\"amlich der Approximation von
L\"osungsoperatoren partieller Differentialgleichungen (PDEs) der Form
$\mathcal{G}^\dagger\colon \mathcal{X} \to \mathcal{Y}$, die von einem
Eingaberaum (z.\,B.\ Anfangsbedingungen, Koeffizienten) in einen
L\"osungsraum abbilden.

Im Gegensatz dazu arbeitet der weit verbreitete \emph{Fourier Neural Operator}
(FNO) im Spektralraum und nutzt globale Faltungen in der Fourierbasis.
Obwohl der FNO starke theoretische Garantien besitzt, weist er in der Praxis
Schwachstellen auf: mangelnde Translationsäquivarianz, Aliasing-Artefakte bei
h\"oheren Frequenzen sowie Schwierigkeiten bei der Darstellung lokaler Strukturen.

Die vorliegende Arbeit reagiert auf diesen Zustand mit der Entwicklung des
\textbf{Convolutional Neural Operator (CNO)}, der die Ausdruckskraft und
rechentechnische Effizienz von CNNs mit rigorosen theoretischen Garantien
f\"ur das Operator"=Lernen verbindet.

% =========================================================
\subsection*{Kernbeitrag: Der Convolutional Neural Operator (CNO)}
% =========================================================

Der CNO ist eine modifizierte U-Net-Architektur, die so konstruiert ist, dass
sie \emph{bandlimitierte Funktionen} als Ein- und Ausgaben verarbeitet. Die
zentrale theoretische Grundlage ist das Konzept der
\textbf{kontinuierlich-diskreten \"Aquivalenz} (Continuous-Discrete Equivalence,
CDE): Die Architektur gew\"ahrleistet, dass ihre diskrete Computerimplementierung
exakt der zugrundeliegenden kontinuierlichen Operatorabbildung entspricht.

Der CNO ist formell definiert als eine Komposition
\[
  \mathcal{G}\colon u \;\mapsto\; P(u) = v_0 \mapsto v_1 \mapsto \cdots
  \mapsto v_L \mapsto Q(v_L),
\]
wobei jede Zwischenschicht die Form $v_{\ell+1} = P_\ell \circ \Sigma_\ell \circ
K_\ell(v_\ell)$ hat. Die drei Grundbausteine unterscheiden sich fundamental von
denen eines standard CNNs:

\begin{itemize}
  \item \textbf{Faltungsoperator $K_\ell$:} Faltung mit kompakten Kernen im
    \emph{Ortsraum} auf einem gleichm\"a\ss igen Gitter, das die
    Nyquist-Bedingung erf\"ullt. Dies steht im direkten Kontrast zum FNO, der
    im Fourier-Raum faltet und damit \emph{global} (nicht lokal) operiert.

  \item \textbf{Up- und Downsampling ($U_{w,\tilde{w}}$, $D_{w,\tilde{w}}$):}
    Anstelle von bilinearer Interpolation oder Average Pooling werden
    \emph{Sinc-Filter} verwendet, die die Bandlimitierung der verarbeiteten
    Funktion strikt einhalten. In der Praxis werden fensterbasierte Sinc-Filter
    (Windowed-Sinc-Filter) mit $N_\text{tap} = 12$ Taps eingesetzt, um das
    Gibbs-Ph\"anomen zu begrenzen.

  \item \textbf{Alias-freie Aktivierungsschicht $\Sigma_{w,\tilde{w}}$:}
    Statt der punktweisen Anwendung einer Aktivierungsfunktion -- die zwangsl\"aufig
    Frequenzinhalte \"uber die Bandgrenze hinaus erzeugt -- wird zun\"achst auf eine
    h\"ohere Bandbreite $\tilde{w} > w$ hochgesampelt, dann die Aktivierung
    angewendet, und anschlie\ss end auf die Ursprungsbandbreite $w$
    zur\"uckgesampelt. Mit \"Ubersamplingfaktor $N_\sigma = 2$ ist dies in der
    Praxis ausreichend und rechentechnisch effizient.
\end{itemize}

Die Gesamtarchitektur folgt dem U-Net-Prinzip: Ein Encoder komprimiert die
r\"aumliche Aufl\"osung bei gleichzeitiger Vergr\"o\ss erung der Kanalbreite,
ein Decoder expandiert die Aufl\"osung zur\"uck, und Skip-Connections
zwischen symmetrischen Ebenen \"ubertragen Hochfrequenzanteile direkt. Zus\"atzlich
werden ResNet-Bl\"ocke zwischen Encoder- und Decoder-Ebenen eingef\"ugt, die sich
in Ablationsstudien als essenziell f\"ur gute Generalisierungsleistung erweisen.

% =========================================================
\subsection*{Theoretische Garantien}
% =========================================================

Die Autoren beweisen zwei zentrale theoretische Resultate:

\begin{enumerate}
  \item \textbf{Repr\"asentations\"aquivalenz (Proposition 2.1):} Der CNO ist
    ein \emph{Representation Equivalent Neural Operator} (ReNO) im Sinne von
    Bartolucci et al.\ (2023). Das bedeutet, er erf\"ullt die
    CDE-Eigenschaft exakt und ist damit resolutionsinvariant: Ein auf einer
    bestimmten Gitteraufl\"osung trainiertes Modell kann ohne Neutraining auf
    beliebige andere Aufl\"osungen angewendet werden, ohne Aliasing-Fehler
    einzuf\"uhren.

  \item \textbf{Universeller Approximationssatz (Theorem 3.1):} F\"ur jeden
    stetigen PDE-L\"osungsoperator $\mathcal{G}^\dagger$ und jedes
    $\varepsilon > 0$ existiert ein CNO $\mathcal{G}$, sodass
    $\|\mathcal{G}^\dagger(a) - \mathcal{G}(a)\|_{L^p} < \varepsilon$ f\"ur
    alle zul\"assigen Eingaben $a$. Der Beweis nutzt das universelle
    Approximationstheorem f\"ur feedforward-Netze sowie trigonometrische
    Polynominterpolation und gilt f\"ur eine breite Klasse von PDEs,
    einschlie\ss lich elliptischer, hyperbolischer und parabolischer Gleichungen.
\end{enumerate}

% =========================================================
\subsection*{Experimentelle Ergebnisse}
% =========================================================

Der CNO wird auf acht repr\"asentativen PDE-Benchmarks (RPB) evaluiert, die
bewusst gew\"ahlt wurden, um Probleme mit mehreren r\"aumlichen Skalen zu
umfassen -- Probleme also, bei denen klassische numerische L\"oser auf feinen
Gittern kostspielig sind und ML-Surrogate einen echten Vorteil bieten.
Die wichtigsten Resultate im Vergleich zu den Baselines (U-Net, FNO,
DeepONet, ResNet, FFNN, Galerkin Transformer) lassen sich wie folgt
zusammenfassen:

\begin{itemize}
  \item \textbf{In-distribution Genauigkeit:} Der CNO \"ubertrifft alle
    Baselines auf sieben von acht Benchmarks. Besonders ausgepr\"agt ist der
    Vorsprung gegen\"uber dem FNO bei der Poisson-Gleichung (Faktor~20) und
    den Navier-Stokes-Gleichungen.

  \item \textbf{Out-of-distribution Generalisierung:} Auf Transportprobleme
    (Smooth und Discontinuous Transport) generalisiert der FNO schlecht, da er
    keine Translationsäquivarianz besitzt. CNO und U-Net, die r\"aumlich
    lokal operieren, generalisieren deutlich robuster.

  \item \textbf{Resolutionsinvarianz:} W\"ahrend der Testfehler des FNO bei
    Ver\"anderung der Gitteraufl\"osung um bis zu 25\,\% schwankt und der
    Fehler des Standard-U-Net sogar um den Faktor~3 ansteigt, bleibt der
    CNO-Fehler \"uber alle getesteten Aufl\"osungen nahezu konstant.

  \item \textbf{Recheneffizienz und Datenskalierung:} F\"ur denselben
    Modellumfang liefert der CNO kleinere Validierungsfehler als der FNO und
    ist dabei fast doppelt so schnell zu trainieren. Die Datenskalierung folgt
    einem Potenzgesetz mit Exponent $r = 0.37$ (CNO) gegen\"uber $r = 0.28$
    (FNO), was bedeutet, dass der CNO etwa viermal weniger Trainingsdaten
    ben\"otigt, um denselben Fehler von 1\,\% zu erreichen.

  \item \textbf{Stabilit\"at:} \"Uber zehn verschiedene Zufallsinitialisierungen
    zeigt der CNO sehr geringe Standardabweichungen des Testfehlers, was auf
    ein robustes Trainingsverhalten hindeutet.
\end{itemize}

% =========================================================
\subsection*{Bezug zur Masterarbeit}
% =========================================================

Die CNO-Arbeit ist in mehrfacher Hinsicht von hoher Relevanz f\"ur die
vorliegende Masterarbeit, die U-Net und FNO als Surrogatmodelle f\"ur
Str\"omungsfelder in Systemen mit sinkenden Kristallen im Stokes-Regime
vergleicht.

\subsubsection*{1.\ Theoretische Einordnung des U-Net/FNO-Vergleichs}

Ein zentrales Ergebnis der CNO-Arbeit ist, dass das standard U-Net -- trotz
seiner guten empirischen Leistung -- kein ReNO ist und daher keine CDE
erf\"ullt. Es ist also streng genommen kein konsistenter Operatorlerner,
sondern ein diskretisierungsabh\"angiges Modell. Der FNO erfüllt die CDE
in der Theorie, zeigt aber in der Praxis aufgrund von Implementierungsdetails
(z.\,B.\ Gitterinterpolation) ebenfalls Aufl\"osungsabh\"angigkeiten.

F\"ur die Masterarbeit bedeutet dies: Die beobachteten Leistungsunterschiede
zwischen U-Net und FNO lassen sich nicht allein auf Hyperparameter
zur\"uckf\"uhren, sondern reflektieren fundamentale Unterschiede in den
\emph{Induktiven Biases} beider Architekturen. Das U-Net bevorzugt lokal
konsistente, r\"aumlich hierarchische Repr\"asentationen, der FNO globale
spektrale Abbildungen. Im Stokes-Regime, das durch elliptische Fernwirkung
gepr\"agt ist, lautet die relevante Frage: Welcher Bias passt besser zum
physikalischen Problem?

\subsubsection*{2.\ Lokale vs.\ globale Verarbeitung bei Kristallgeometrien}

Die CNO-Arbeit zeigt, dass lokal operierende Architekturen (U-Net, CNO)
bei Transportproblemen und Problemen mit lokalen Strukturen -- Diskontinuit\"aten,
scharfen Grenzfl\"achen -- klare Vorteile gegen\"uber dem FNO haben. Dieses
Ergebnis ist f\"ur die Masterarbeit von direkter Bedeutung:
\begin{itemize}
  \item Die Str\"omungsfelder um Kristallgrenzen herum weisen steile
    Geschwindigkeitsgradienten auf, die hohe Frequenzinhalte erfordern und
    vom FNO mit begrenzten Fourier-Moden nur unzureichend aufgel\"ost werden
    k\"onnen (vgl.\ Spektralverzerrungen in Figure~2 der CNO-Arbeit).
  \item Das U-Net hingegen \"ubertr\"agt Hochfrequenzanteile \"uber
    Skip-Connections direkt zwischen Encoder- und Decoder-Ebenen, was es
    strukturell besser f\"ur die Darstellung scharfer Kristallgrenzen geeignet
    macht.
  \item Gleichzeitig profitiert der FNO von der globalen elliptischen Natur
    der Stokes-Gleichungen ($\nabla \cdot \mathbf{v} = 0$, Inkompressibilit\"at
    als Fernwirkung), die spektrale Repr\"asentationen beg\"unstigt.
\end{itemize}
Diese komplementären St\"arken motivieren direkt den systematischen
Vergleich in der Masterarbeit.

\subsubsection*{3.\ Alias-freie Aktivierungen und der Kontinuit\"atsverlust}

Ein kritischer Befund der CNO-Arbeit ist, dass naiv eingesetzte Aktivierungsfunktionen -- auch in einer U-Net-Architektur -- Aliasing-Fehler
erzeugen, die die Str\"omungsrekonstruktion bei hohen Frequenzen verschlechtern.
Dies ist besonders relevant f\"ur die Stream-Funktion-Repr\"asentation in der
Masterarbeit: Die Divergenzfreiheit ($\nabla \cdot \mathbf{v} = 0$) wird durch
die Stromfunktionformulierung zwar mathematisch garantiert, aber die Aktivierungsschichten k\"onnen bei der Rekonstruktion von $\psi$ nichtphysikalische
Hochfrequenzinhalte einf\"uhren. Der CNO-Ansatz -- Upsampling vor der
Aktivierung, Downsampling danach -- k\"onnte als zus\"atzliche Regularisierungsma\ss nahme in k\"unftigen Modellversionen der Masterarbeit in Betracht
gezogen werden.

\subsubsection*{4.\ Generalisierung \"uber Konfigurationen}

Die CNO-Arbeit untersucht out-of-distribution Generalisierung (z.\,B.\ training
mit $K{=}16$ Fourier-Moden, Test mit $K{=}20$) und zeigt, dass der CNO dabei
deutlich robuster als der FNO ist. Die Kernfrage der Masterarbeit -- ob ein auf
Konfigurationen mit bis zu $N$ Kristallen trainiertes Modell auf Konfigurationen
mit $N + m$ Kristallen generalisiert -- ist eine direkte Analogie zu diesem
out-of-distribution Problem. Die Erkenntnisse der CNO-Arbeit legen nahe, dass
Translationsäquivarianz (wie sie sowohl CNO als auch U-Net besitzen) ein
wichtiger Faktor f\"ur solche Generalisierungsleistungen ist: Ein Modell, das
gelernt hat, die Str\"omungswechselwirkung von zwei Kristallen in beliebigen
r\"aumlichen Konfigurationen zu beschreiben, kann dieses Wissen leichter
auf drei oder mehr Kristalle \"ubertragen, wenn die Verarbeitung
positionsinvariant ist.

\subsubsection*{5.\ CNO als m\"ogliche Hybridarchitektur f\"ur zuk\"unftige Arbeit}

Die CNO-Architektur stellt konzeptionell eine direkte Synthese der beiden
in der Masterarbeit verglichenen Ans\"atze dar: Sie nutzt die r\"aumlich-lokale
Verarbeitung des U-Net (Ortsfaltungen, hierarchische Encoder-Decoder-Struktur)
und erg\"anzt sie um theoretisch fundierte Aliasing-Freiheit, die im FNO
durch die Fourier-Repr\"asentation erreicht wird. Sollte keiner der beiden
Ans\"atze -- reines U-Net oder reines FNO -- die Zielschwelle von
MAE $< 0.05$ f\"ur eine exzellente Bewertung erreichen, bietet der CNO
eine nat\"urliche Erweiterung. Die Open-Source-Implementierung der Autoren
(PyTorch, GitHub) erm\"oglicht eine direkte Adaption auf die
Julia/Flux.jl-Umgebung der Masterarbeit, wenngleich eine Portierung der
Sinc-Filter-Implementierung (CUDA-basiert in der Originalarbeit) zus\"atzlichen
Aufwand erfordert.

% =========================================================
\subsection*{Kritische Einsch\"atzung}
% =========================================================

Die CNO-Arbeit pr\"asentiert \"uberzeugende Ergebnisse, weist jedoch einige
Einschr\"ankungen auf, die f\"ur die Masterarbeit relevant sind:

\begin{itemize}
  \item Alle Benchmarks sind zweidimensional; die Erweiterung auf 3D wird als
    konzeptionell straightforward, aber rechenaufw\"andig beschrieben. F\"ur
    die 3D-Komponente der Masterarbeit ($32^3$ Aufl\"osung) k\"onnte der
    Rechenaufwand der Sinc-Filter prohibitiv sein.
  \item Die Benchmarks umfassen ausschlie\ss lich kartesische Dom\"anen. Nicht-kartesische
    Geometrien (z.\,B.\ k\"unstliche Randgeometrien um Kristalle) erfordern
    weitere Transformationen.
  \item Die Trainingsmengen (512--1024 Samples) liegen in einer vergleichbaren
    Gr\"o\ss enordnung wie die Masterarbeit (200--300 Samples), sodass die
    Datenskalierungsergebnisse (Potenzgesetz-Exponent $r = 0.37$) direkt
    \"ubertragbar sind.
\end{itemize}

\end{document}